% ATTEMPT_STATUS: unresolved
% RESEARCH_GENERATED: 2026-10-01; GPT-6 Astra Ultra; individual mathematical attempt
% TOP_PROBLEM: {"id": 164, "title": "Optimal Sidon Bases", "category_id": 2, "category": {"id": 2, "name": "combinatorics", "display_name": "Combinatorics", "description": "Counting problems, graph theory, discrete structures.", "slug": "combinatorics", "order_index": 2, "created_at": "2026-07-31T15:26:25.671Z"}, "status": "open"}
% FIRST_POSTED: 2026-10-01
% Imported from: unsolvedmath_additional_500.tex; UnsolvedMath ID 164
% !TeX program = lualatex
% Compile twice: lualatex 164.tex
% Self-contained: no external bibliography, images, or \input files.
% Numeric section labels and visible numbers are original UnsolvedMath IDs.
\documentclass[11pt,a4paper]{article}
\usepackage[margin=24mm,headheight=14pt]{geometry}
\usepackage{fontspec}
\setmainfont{Latin Modern Roman}
\setsansfont{Latin Modern Sans}
\setmonofont{Latin Modern Mono}
\usepackage{amsmath,amssymb,mathtools}
\usepackage{unicode-math}
\setmathfont{Latin Modern Math}
\usepackage{microtype}
\usepackage{enumitem,xurl,needspace}
\usepackage[dvipsnames]{xcolor}
\usepackage{hyperref}
\hypersetup{unicode=true,colorlinks=true,linkcolor=MidnightBlue,urlcolor=MidnightBlue,
 pdftitle={UnsolvedMath Problem 164},pdfauthor={Source-annotated compilation},bookmarksnumbered=true}
\usepackage{bookmark,tocloft,titlesec,fancyhdr}
\setlength{\cftsecnumwidth}{4.2em}
\cftsetpnumwidth{2.5em}
\cftsetrmarg{4em}
\titleformat{\section}{\Large\bfseries\raggedright}{\thesection}{0.7em}{}
\pagestyle{fancy}\fancyhf{}
\fancyhead[L]{\small\sffamily UnsolvedMath research attempt}
\fancyhead[R]{\small Original catalogue IDs}
\fancyfoot[C]{\thepage}
\setlength{\parindent}{0pt}
\setlength{\parskip}{5pt plus 1pt}
\setlength{\emergencystretch}{3em}
\setcounter{tocdepth}{1}\setcounter{secnumdepth}{1}
\setlist[itemize]{leftmargin=3em,itemsep=4pt,topsep=4pt,parsep=0pt}
\allowdisplaybreaks
\newcommand{\N}{\mathbb{N}}
\newcommand{\Z}{\mathbb{Z}}
\newcommand{\Q}{\mathbb{Q}}
\newcommand{\R}{\mathbb{R}}
\newcommand{\C}{\mathbb{C}}
\newcommand{\F}{\mathbb{F}}
\newcommand{\A}{\mathbb{A}}
\newcommand{\PP}{\mathbb{P}}
\newcommand{\E}{\mathbb{E}}
\newcommand{\eps}{\varepsilon}
\newcommand{\dd}{\,\mathrm d}
\newcommand{\sref}[2]{\hyperlink{source:#1:#2}{[S#2]}}
\newif\ifshowcatalogue
\showcataloguetrue
\newcommand{\cataloguescope}[1]{\ifshowcatalogue
 \paragraph{Statement recorded in the supplied source.}
 #1\fi}
\begin{document}
% UnsolvedMath ID 164; source catalogue code GREEN-076

\setcounter{section}{163}

\section{Near-Optimal Twofold Bases of Cyclic Groups}\label{164}

{\small\sffamily UnsolvedMath ID 164 · Combinatorics}

\subsection{Definitions and mathematical statement}

\paragraph{Shared notation.}Unless a section says otherwise, $\N=\{1,2,\ldots\}$,
$\Z$ is the ring of integers, $\Q,\R,\C$ are the rational, real, and complex
fields, $[N]=\{1,\ldots,\lfloor N\rfloor\}$, and $\log$ is the natural
logarithm. For real functions of a parameter tending to infinity,
$f\ll g$ and $f=O(g)$ mean $|f|\leq Cg$ eventually for some fixed $C$;
$f\gg g$ reverses this comparison; $f=o(g)$ means $f/g\to0$;
$f\sim g$ means $f/g\to1$; and $f\asymp g$ means both order bounds.
A subscript on $O$ or $\ll$ specifies allowed constant dependence.
The expression $n^{o(1)}$ in an upper bound means growth smaller than
$n^\epsilon$ for each fixed $\epsilon>0$. Local definitions govern any
exception, finite-field convention, or use of the same symbol for another
object. An asymptotic claim is not automatically an effective algorithm.



A twofold additive basis of $\Z/q\Z$ is a set $A$ with $A+A=\Z/q\Z$. There are at most $|A|(|A|+1)/2$ unordered sums, giving the counting scale $\sqrt{2q}$. The source asks for an unbounded sequence of moduli $q$ admitting bases with $|A|/\sqrt q\to\sqrt2$. Uniqueness of representations is not required. \sref{164}{1} \sref{164}{2}

\paragraph{Statement recorded in the supplied source.}
Are there infinitely many $q$ for which there is a set $A \subset \mathbb{Z}/q\mathbb{Z}$ with $|A| = (\sqrt{2} + o(1))q^{1/2}$ and $A + A = \mathbb{Z}/q\mathbb{Z}$? \sref{164}{1}

\subsection{Short English statement}

Can cyclic groups of arbitrarily large orders have two-summand additive bases that asymptotically attain the elementary counting lower bound?

\subsection{Research attempt}

\paragraph{Provenance and scope.}
This individual attempt was generated by GPT-6 Astra Ultra on 1 October 2026. The method was to derive explicit partial results and test the first proposed extension adversarially. The original statement and sources below are retained. No claim of mathematical novelty or complete solution is made.

\paragraph{Formal target and counting convention.}
We seek cyclic groups of unbounded orders $q$ with a two-summand basis $A$ of size $(\sqrt2+o(1))\sqrt q$. Repetitions are allowed, but commutativity means unordered pairs are the relevant first count. The primary problem list distinguishes this cyclic question from analogous geometric direction problems. We establish an explicit constant-two construction, rule out exact equality in the counting bound beyond the smallest cases, and prove that any asymptotically optimal family must develop unbounded representation multiplicities.

\paragraph{The elementary lower bound and a complete construction.}
If $m=|A|$, then $m(m+1)/2$ unordered pairs must cover $q$ residues. Thus
\[
 m\geq\frac{\sqrt{8q+1}-1}{2}=\sqrt{2q}+O(1).
\]
For $q=t^2$, take
\[
 A=\{0,1,\ldots,t-1\}\cup\{0,t,2t,\ldots,(t-1)t\}.
\]
It has $2t-1$ elements. Every residue has a representative $jt+i$ with $0\leq i,j<t$, the sum of one element from each displayed part. Hence $A+A=\Z/t^2\Z$. This proves a constant-two upper construction along infinitely many orders. The shared zero is counted once, and the proof does not require primality.

\paragraph{Exact perfect unordered coverage is impossible for large sets.}
Suppose equality $q=m(m+1)/2$ holds and every residue has exactly one unordered representation. If $a-b=c-d$ with $a,b,c,d\in A$ and $a\ne b$, then $a+d=c+b$. Unique unordered representation forces either $a=c,d=b$, or $a=b,d=c$. The second alternative contradicts $a\ne b$. Thus all $m(m-1)$ ordered nonzero differences are distinct, so $m(m-1)\leq q-1$. Combining the two formulas gives
\[
 m^2-3m+2\leq0,
\]
forcing $m\leq2$. The case $m=2,q=3$, with $A=\{0,1\}$, indeed attains exact equality. Therefore the desired asymptotic cannot arise from genuinely perfect unique-sum designs of growing size. Small relative waste, rather than zero waste, is indispensable.

\paragraph{A quantitative collision constraint on near-optimal bases.}
Let $u(s)$ be the number of unordered representations of $s$, and let $R=\max_su(s)$. Since $A$ is a basis, $u(s)\geq1$. Define its excess-pair budget
\[
 D=\sum_s(u(s)-1)=\frac{m(m+1)}2-q.
\]
The ordered additive energy $E(A)$ can be counted by differences. The zero difference contributes $m^2$, and the remaining difference multiplicities have total $m(m-1)$. Cauchy--Schwarz over at most $q-1$ nonzero residues gives
\[
 E(A)\geq m^2+\frac{m^2(m-1)^2}{q-1}.
\]
On the sum side, every unordered pair gives at most two ordered ones, so
\[
 E(A)\leq4\sum_su(s)^2
 \leq4q+4(R+1)D.
\]
The last step uses $u^2-1=(u-1)(u+1)\leq(R+1)(u-1)$. Therefore
\[
 (R+1)D\geq\frac14\left(m^2+
 \frac{m^2(m-1)^2}{q-1}-4q\right).
\]
If $m=(\sqrt2+o(1))\sqrt q$, the right-hand side is $(1/2+o(1))q$, while $D=o(q)$. It follows that $R\to\infty$. Thus an asymptotically optimal basis cannot have all unordered representation counts bounded by one fixed constant, even though all but $o(q)$ residues must have a unique representation because their number of exceptions is at most $D$.

\paragraph{Meaning of the obstruction.}
These two conclusions are compatible: a vanishing proportion of residues may carry increasingly large multiplicities, supplying the energy forced by the difference count. A strategy that tries to construct an almost-Sidon basis with a uniform bounded representation cap is therefore ruled out by the displayed inequality. It would be incorrect to conclude that near-optimal bases themselves are impossible; the argument leaves exactly this concentration mechanism available.

\paragraph{Verification and remaining gap.}
The difference argument respects ordered differences and unordered sums separately, and the energy upper bound remains valid for even $q$ without assuming doubled elements are distinct. The construction works for every positive integer $t$, with the asymptotic question using $t\to\infty$. The proved results delimit the structure of any optimal family but do not construct one. The remaining task is to realise a cyclic basis with small excess budget and sufficiently concentrated representation collisions, or to prove an additional obstruction preventing that combination.

\subsection{Sources}

{\small

\begin{itemize}

\item[\hypertarget{source:164:1}{S1}] ULAM.AI, \emph{UnsolvedMath}, supplied \texttt{problems.json}, record 164 (GREEN-076), statement and background. The embedded literature-review date is 2026-08-17; that is the archive’s review date, not a fresh certification. Dataset landing page: \url{https://huggingface.co/datasets/ulamai/UnsolvedMath}.

\item[\hypertarget{source:164:2}{S2}] Ben Green, 100 Open Problems, Problem 33, most recent update December 2025. \url{https://people.maths.ox.ac.uk/greenbj/papers/open-problems.pdf}. The author-hosted document was inspected on 27 September 2026; the archive’s local code is not a problem-number locator in this paper.

\end{itemize}

}

\label{end:164}
\end{document}
